\subsection{Sistemas LTI}

\begin{frame}
\frametitle{Sistemas LTI}

\vspace{-0,3cm}

  \begin{figure}
    \centering
    \small
    \begin{tikzpicture}
      \node[bloque,minimum width=15mm] (T) {LTI};
      \node (x) [left=of T] {};
      \node (y) [right=of T] {};
      \draw[arrow,->] (x.east) -- (T.west) node [midway,above] {$x(\cdot)$};
      \draw[arrow,->] (T.east) -- (y.west) node [midway,above] {$y(\cdot)$};
    \end{tikzpicture}
  \end{figure}

  \begin{itemize}
    \item <1-> Constituyen una clase de sistemas muy importantes en procesamiento de señales, control de sistemas, teoría de la comunicación.

    \item Características más relevantes:

    \begin{itemize}

      \item <1-> La respuesta de un sistema LTI a una entrada puede obtenerse como la \alert{convolución} de la entrada con la respuesta del sistema a la entrada impulso unitario.

      \item <1-> Las \alert{transformadas} de Laplace, Z, Fourier \alert{de la convolución es el producto de las transformadas}. 

      \item Por lo tanto, podemos calcular la respuesta de los sistemas LTI ante una entrada a partir del producto de la transformada de la entrada y de la transformada del sistema. 

    \end{itemize}

  \end{itemize}

  \vspace{-0.4cm}

\end{frame}

%%%%%%%%%%%%%%%%%%%%%%%%%%%%%%%%%%%%%%%%%%%%


\begin{frame}
\frametitle{Respuesta temporal sistema LTI}

\vspace{-0,3cm}


\begin{itemize}
\item <1-> \alert{Sistemas Continuos}:

\begin{itemize}
\item Respuesta al impulso (función ponderatriz):   $\bm{T}[\delta(t)]=h(t)$
	
	\vspace{0,1cm}
	
\item Respuesta a una entrada arbitraria:

\vspace{0,2cm}

\footnotesize $\begin{array}{ll}\color{red}\!\!\bm{T}[x(t)]\color{black}\!\!\!\!\!\!&=\bm{T}[(x\!*\!\delta)(t)]\!=\!\bm{T}[\int_{-\infty}^{\infty} \! x(\tau)\delta({t\!-\!\tau})d\tau] \!=\!\int_{-\infty}^{\infty}\bm{T}[x(\tau)\delta({t\!-\!\tau})]d\tau \!=\\ [0.2cm]  & =\int_{-\infty}^{\infty} x(\tau) \bm{T}[\delta({t\!-\!\tau})]d\tau  =\int_{-\infty}^{\infty} x(\tau)  h(t\!-\!\tau) d\tau \color{red}= (x*h)(t)\color{black} \end{array}$


\end{itemize}


\item <2-> \alert{Sistemas Discretos}:

\begin{itemize}
\item Respuesta al impulso (función ponderatriz):   $\bm{T}[\delta(n)]=h(n)$
	
	\vspace{0,1cm}
	
\item  Respuesta a una entrada arbitraria:

\vspace{0,1cm}

\footnotesize $\begin{array}{ll}\color{red}\!\!\bm{T}[x(n)]\color{black}\!\!\!\!\!\!&=\!\bm{T}[(x\!*\!\delta)(n)]\!=\!\bm{T}[\sum_{k=-\infty}^{\infty} \! x(k)\delta({n\!-\!k})] \!=\!\!\sum_{k=-\infty}^{\infty}\! \!\bm{T}[x(k)\delta({n\!-\!k})]\!=\\ [0.2cm]  & =\!\sum_{k=-\infty}^{\infty} x(k) \bm{T}[\delta({n\!-\!k})]  \!=\!\sum_{k=-\infty}^{\infty} x(k)  h(n\!-\!k) \color{red} \!=\! (x\!*\!h)(n)\color{black} \end{array}$

\vspace{0,2cm}
\end{itemize}


\item <3-> \small La respuesta de un  sistema LTI a cualquier entrada $x(\cdot)$ se puede calcular como la convolución de la entrada y la respuesta del sistema a la entrada impulso $h(\cdot)$.

\vspace{0,1cm}

\hspace{2cm} \alert{$h(\cdot)=\bm{T}[\delta(\cdot)]  \rightarrow  \bm{T}[x(\cdot)]=(x* h)(\cdot)$}


\end{itemize}


\end{frame}

%%%%%%%%%%%%%%%%%%%%%%%%%%%%%%%%%%%%%%%%%%%%


\begin{frame}
\frametitle{Sistemas LTI continuos I}

\vspace{-0,3cm}


\begin{itemize}
\item <1-> \alert{Transformada de Laplace (TL)}:

\begin{itemize}
\item Bilateral:    $\mathcal{L}{[x(t)]}\!=\!X(s)\!=\! \int_{-\infty}^{\infty} \! \!x(t) e^{-t s} dt \:$ con   $s \in \mathbb{C}$


\vspace{0.1cm}

\item \color{olivegreen} Unilateral:    $\mathcal{L}{[x(t)]}\!=\!X(s)\!= \!\int_{0}^{\infty} \! \! x(t) e^{-t s}dt  \:$ con   $s \in \mathbb{C}$, $x(t)$ causal \color{black}

\vspace{-0.4cm}

\end{itemize}

\item <2-> \alert{Propiedades} (similares, se diferencian en las condiciones iniciales que aparecen en las derivadas):

\begin{itemize}
\item Linealidad : $\mathcal{L}[\alpha x_1(t)+\beta  x_2(t)]=\alpha \mathcal{L}[x_1(t)]+\beta \mathcal{L}[x_2(t)]$
\vspace{0.2cm}
\item Desplazamiento temporal: $\mathcal{L}[x(t-t_0)]=e^{-st_0} \mathcal{L}[x(t)]$
\vspace{0.2cm}
\item Derivada: $\mathcal{L}[\frac {dx(t)}{dt}]=s \mathcal{L}[x(t)]\color{olivegreen}-x(0)\color{black}$
\vspace{0.2cm}
\item Convolución: $\mathcal{L}[{(x*y)(t)}]= \mathcal{L}[x(t)]\mathcal{L}[y(t)]$
\end{itemize}

\item <3-> \alert{TL de la respuesta del sistema}:

\vspace{0.2cm}

\footnotesize $\mathcal{L}[ y(t)]\!=\!\mathcal{L}[ \bm T[{(x)(t)}]\!=\!\mathcal{L}[{(x*h)(t)}]\!=\!\mathcal{L}[{x(t)}]\mathcal{L}[{h(t)}]   \leftrightarrow     Y(s)\!=\!H(s)X(s)$  

\vspace{0.2cm}

La respuesta de un sistema LTI continuo a una entrada se puede calcular a partir de la TL inversa del producto de la TL de la señal y de la TL de la respuesta del sistema a la entrada impulso. 
\end{itemize}

%%%%%%%%%%%%%%%%%%%%%%%%%%%%%%%%%%%%%%%%%%%%%%%%%%%%%%%%%%%%%%%%%%%%%%%%

\end{frame}

\begin{frame}
\frametitle{Sistemas LTI continuos II}

\vspace{-0,3cm}


\begin{itemize}

\item <1-> \alert{TL de la respuesta del sistema}:

\vspace{0.2cm}

\hspace{1cm}{\footnotesize $\mathcal{L}[ y(t)]\!=\!\mathcal{L}[{x(t)}]\mathcal{L}[{h(t)}]   \leftrightarrow     Y(s)\!=\!H(s)X(s)$  }

\vspace{0.2cm}

\item <2-> \alert{Ecuación diferencial de un sistema LTI continuo:}

\vspace{0.2cm}


\hspace{0.5cm} {\footnotesize $\sum_{k=0}^N {a_k \frac{d^k y(t)}{dt}}=\sum_{k=0}^M {b_k \frac{d^k x(t)}{dt}}$   $\buildrel TL \over  \longrightarrow$   $\sum_{k=0}^N a_k s^k Y(s)=\sum_{k=0}^M b_k s^k X(s)  \rightarrow $

\vspace{0.2cm}

\hspace{0.5cm} $Y(s)=\frac{\sum_{k=0}^M b_k s^k}{\sum_{k=0}^N a_k s^k}X(s)=H(s)X(s)$}

\vspace{0.2cm}

\item $H(s)$, la TL de la respuesta del sistema LTI a la función impulso, es la función de transferencia del sistema.

\vspace{0.2cm}

\item <3-> \alert{Condiciones adicionales:}

\begin{itemize}
\item Sistema dinámico $N \ge 1$
\item Sistema causal: 
\begin{itemize} 
\item $M \le N$
\item $h(t)=T[\delta(t)]$,  $h(t)=0 \: \: t<0$
\end{itemize}
\item Estable: Polos (ceros del denominador de $H(s)$) tienen que tener la parte real negativa.
\end{itemize}


\end{itemize}


\end{frame}



%%%%%%%%%%%%%%%%%%%%%%%

\begin{frame}\frametitle{Sistemas LTI continuos III}\small
  \begin{figure}\centering\small
    \begin{tikzpicture}
      \node[bloque,minimum width=15mm] (T) {LTI};
      \node (x) [left=of T] {};
      \node (y) [right=of T] {};
      \draw[arrow,->] (x.east) -- (T.west) node [midway,above] {$x(\cdot)$};
      \draw[arrow,->] (T.east) -- (y.west) node [midway,above] {$y(\cdot)$};
    \end{tikzpicture}
  \end{figure}
  \begin{itemize}\footnotesize
    \item Transformada  de Fourier respuesta $u(t)=\delta(t)$: $H(j\omega)=H(s=j\omega)$
    \item Respuesta permanente LTI estable a $u(t)=e^{j\omega t}$
    
    \footnotesize
    $\alert{y(t)=}\bm{T}[e^{j\omega t}]=\int_{-\infty}^\infty h(\tau)e^{j(\omega(t-\tau))}d\tau=\int_{-\infty}^\infty h(\tau)e^{j\omega t}e^{-j\omega\tau}d\tau=$

    \hspace{0.6cm}	 $=e^{j\omega t} \int_{-\infty}^\infty h(\tau)e^{-j\omega\tau}d\tau=\alert{H(j\omega)e^{j\omega t}}=|H(j\omega)|e^{j(\omega t+arg(H(j\omega)))}$
    \item Respuesta permanente LTI estable a $u(t)=sin(\omega t)$:

    $\alert{y(t)=|H(j\omega)|sen(\omega t+arg(H(j\omega)))}$
    \item Respuesta permanente LTI estable a $u(t)$ periodica ($T$, $\omega=\frac{2\pi}{T}$):

    $\alert{y(t)=\sum_{k=-\infty}^{\infty} c_k |H(jk\omega)|e^{j(\omega kt+arg(H(jk\omega)))}}$
    \item Respuesta permanente  LTI estable a $u(t)$ gen\'{e}rica: $Y(j\omega)=H(j\omega)U(j\omega)$
  \end{itemize}
\end{frame}

%%%%%%%%%%%%%%%%%%%%%%%%%%%%%%%%%%%%%%%%%%%%


\begin{frame}
\frametitle{Sistemas LTI discretos I}

\vspace{-0,3cm}


\begin{itemize}
\item <1-> \alert{Transformada Z (TZ)}:

\begin{itemize}
\item Bilateral:    $\mathcal{Z}{[x(n)]}\!=\!X(z)\!=\! \sum_{n=-\infty}^{\infty} \! \!x(n) z^{-n}  \:$ con   $z\! \in\! \mathbb{C}$

\vspace{0.1cm}

\item \color{olivegreen} Unilateral:    $\mathcal{Z}{[x(n)]}\!=\!X(z)\!= \!\sum_{n=0}^{\infty} \! \! x(n) z^{-n}  \:$ con   $z\! \in\! \mathbb{C}$, $x(\!n\!)$ causal \color{black}

\vspace{-0.4cm}

\end{itemize}

\item <2-> \alert{Propiedades} (similares, se diferencian en las condiciones iniciales que aparecen en los retardos):

\begin{itemize}
\item Linealidad : $\mathcal{Z}[\alpha x_1(n)+\beta  x_2(n)]=\alpha \mathcal{Z}[x_1(n)]+\beta \mathcal{Z}[x_2(n)]$
\vspace{0.2cm}
\item Desplazamiento temporal: $\mathcal{Z}[x(n-n_0)]=z^{-n_0} \mathcal{Z}[x(n)]$
\vspace{0.2cm}
\item Adelanto: $\mathcal{Z}[ x(n+1)]=z \mathcal{Z}[x(n)]\color{olivegreen}-z x(0)\color{black}$
\vspace{0.2cm}
\item Convolución: $\mathcal{Z}[{(x*y)(n)}]= \mathcal{Z}[x(n)]\mathcal{Z}[y(n)]$
\end{itemize}

\item <3-> \alert{TZ de la respuesta del sistema}:

\vspace{0.2cm}

\footnotesize $\mathcal{Z}[ y(n)]\!=\!\mathcal{Z}[ \bm T[{(x)(n)}]\!=\!\mathcal{Z}[{(x*h)(n)}]\!=\!\mathcal{Z}[{x(n)}]\mathcal{Z}[{h(n)}]   \leftrightarrow     Y(z)\!=\!H(z)X(z)$  

\vspace{0.2cm}

La respuesta de un sistema LTI discreto a una entrada se puede calcular a partir de la TZ inversa del producto de la TZ de la señal y de la TZ de la respuesta del sistema a la entrada impulso. 
\end{itemize}


\end{frame}

%%%%%%%%%%%%%%%%%%%%%%%%%%%%%%%%%%%%%%%%%%%%%%%%%%%%%%%%%%%%%%%%%%%%%%%%


\begin{frame}
\frametitle{Sistemas LTI discretos II}

\vspace{-0,3cm}


\begin{itemize}

\item <1-> \alert{TZ de la respuesta del sistema}:

\vspace{0.2cm}

\hspace{0.5cm}{\footnotesize $\mathcal{Z}[ y(n)]\!=\!\mathcal{Z}[{x(n)}]\mathcal{Z}[{h(n)}]   \leftrightarrow     Y(z)\!=\!H(s)X(z)$  }

\vspace{0.2cm}

\item <2-> \alert{Ecuación en diferencias de un sistema LTI discreto:}

\vspace{0.2cm}


\hspace{-0.2cm} {\footnotesize $\sum_{k=0}^N {a_k {y(n\!+\!k)}}\!=\!\sum_{k=0}^M {b_k  x(n\!+\!k)}$   $\buildrel TZ \over  \longrightarrow$   $\sum_{k=0}^N a_k z^k Y(z)\!=\!\sum_{k=0}^M b_k z^k X(z)  \rightarrow $

\vspace{0.2cm}

\hspace{-0.2cm} $Y(z)=\frac{\sum_{k=0}^M b_k z^k}{\sum_{k=0}^N a_k z^k}X(z)=H(z)X(z)$}

\vspace{0.2cm}

\item $H(z)$, la TZ de la respuesta del sistema LTI a la función impulso, es la función de transferencia del sistema.

\vspace{0.2cm}

\item <3-> \alert{Condiciones adicionales:}

\begin{itemize}
\item Sistema dinámico $N \ge 1$
\item Sistema causal: 
\begin{itemize}
\item $M \le N$
 \item $h(n)=T[\delta(n)]$,  $h(n)=0 \: \: n<0$
\end{itemize}
\item Estable: Polos (ceros del denominador de $H(z)$) tienen que estar dentro del círculo unidad.
\end{itemize}


\end{itemize}
\end{frame}

%%%%%%%%%%%%%%%%%%%%%%%%%%%%%%%%%%%%%%

\begin{frame}
\frametitle{Sistemas LTI discretos III}


\small


\begin{figure}\centering\small
    \begin{tikzpicture}
      \node[bloque,minimum width=15mm] (T) {LTI};
      \node (x) [left=of T] {};
      \node (y) [right=of T] {};
      \draw[arrow,->] (x.east) -- (T.west) node [midway,above] {$x(\cdot)$};
      \draw[arrow,->] (T.east) -- (y.west) node [midway,above] {$y(\cdot)$};
    \end{tikzpicture}
  \end{figure}
  \begin{itemize}\footnotesize
    \item Transformada  de Fourier respuesta $u(n)=\delta(t)$: $H(e^{j\omega})=H(z=e^{j\omega})$
    \item Respuesta permanente LTI estable a $u(n)=e^{j\omega n}$\footnotesize

    $\alert{y(n)=}\bm{T}[e^{j\omega n}]=\sum_{k=-\infty}^\infty h(k)e^{j\omega(n-k)}=\sum_{k=-\infty}^\infty h(k)e^{j\omega n}e^{-j\omega k}=$
    \hspace{0.6cm}	$=e^{j\omega n} \sum_{k=-\infty}^\infty h(k)e^{-j\omega k}=\alert{H(e^{j\omega})e^{j\omega n}}=|H(e^{j\omega})|e^{j(\omega n+arg(H(e^{j\omega})))}$
    \item Respuesta permanente LTI estable a $u(n)=sen(\omega t)$:

    $\alert{y(n)=|H(e^{j\omega})|sen(\omega t+arg(H(e^{j\omega})))}$
    \item Respuesta permanente LTI estable a $u(n)$ periodica ($N$, $\omega=\frac{2\pi}{N}$):

    $\alert{y(n)=\sum_{k=0}^{N-1} c_k |H(e^{jk\omega})|e^{j(\omega kt+arg(H(e^{jk\omega})))}}$
    \item Respuesta permanente  LTI estable a $u(n)$ gen\'{e}rica: $Y(e^{j\omega})=H(e^{j\omega})U(e^{j\omega})$
  \end{itemize}
\end{frame}

%%%%%%%%%%%%%%%%%%%%%%%%%%%%%%%


\begin{frame}
\frametitle{Sistemas LTI discretos IV}

\vspace{-0,0cm}

\hspace{1.3cm} {\footnotesize $y(n)=\bm{T}[x(n)]=(x*h)(n)=(h*x)(n)=\sum_{k=-\infty}^\infty h(k)x(n-k)$

\vspace{0,2cm}

\hspace{1.3cm}  $h(n)=\bm{T}[\delta(n)]$}

\vspace{0,2cm}

\hspace{1.25cm} {\small Sistema causal:} {\footnotesize $h(n)=0$, $n<0$}


\begin{itemize}

\item \small {Tipos de sistemas LTI discretos}:



\begin{itemize}

\item \alert{Sistemas FIR}: Finite Impulse Response

 \vspace{0.1cm}

\hspace{2.2cm} Respuesta finita a la entrada impulso  

 \vspace{0.1cm}

\footnotesize{
\hspace{2.2cm} \alert{$h(n)=0$,  $n \ge L$}

 \vspace{0.2cm}

\hspace{2.2cm}   $y(n)=\sum_{k=0}^{L-1} h(k)x(n-k)$


\hspace{2.2cm}   \color{gray} $M=N$, $N=L-1$, $a_{N-k}\!=\!\left\{\begin{array}{l}  0 \: \: k<N \\ 1 \:\: k=N \end{array}\right.$, $b_{N-k}=h(k)$ \color{black}
}
					

\item  \alert{Sistemas IIR}: Infinite Impulse Response

 \vspace{0.1cm}

\hspace{2.2cm}Respuesta infinita a la entrada impulso  

 \vspace{0.1cm}

{\footnotesize \hspace{2.2cm}$y(n)=\sum_{k=0}^{\infty} h(k)x(n-k)$}


 \vspace{0.1cm}

\hspace{2.2cm}Nos interesan los IIR recursivos


\end{itemize}



\end{itemize}


\vspace{-0,4cm}

\rule{12cm}{0.4pt}

 \color{gray}{\footnotesize $\sum_{k=0}^N {a_k {y(n\!+\!k)}}\!=\!\sum_{k=0}^M {b_k  x(n\!+\!k)}$ $\rightarrow$  $\sum_{k=0}^N {a_{N-k} {y(n\!-\!k)}}\!=\!\sum_{k=N-M}^N {b_{N-k}  x(n\!-\!k)}$}

\vspace{0,2cm}


\end{frame}

%%%%%%%%%%%%%%%%%%%%%%%%%%


\begin{frame}
\frametitle{Sistemas LTI continuos IV}

\small
¿Equivalente LTI continuo a los sistemas FIR e IIR discretos ?

\begin{itemize}
\item No existen sistemas FIR en continuo:
\begin{itemize}
\item \small Sistema FIR discreto:

\vspace{0.1cm}

{\footnotesize $y(n)\!=\!\sum_{k=0}^{L-1} \!h(k)x(n\!-\!k)$ $\rightarrow$   $Y(z)\!=\![h_0\!+\!h_1z^{-1}\!\!+\!...\!+\!h_{L\!-\!1}z^{-L+1}]X(z)$}

\vspace{0.2cm}

Tiene los polos en z=0. 

\vspace{0.2cm}

\item Su equivalente  continuo ($z=e^{Ts}$) tendría que tenerlos en $s=-\infty$. Y eso no es fisicamente realizable.
 

\vspace{0.1cm}

\end{itemize}

\item Los sistemas LTI continuos estables son IIR
\end{itemize}

\end{frame}

